\section{Synthèse}

\begin{cadrebleu}{À retenir}
La statistique descriptive permet de résumer une série de données.

\medskip

\begin{itemize}
    \item Les indicateurs de position situent les données.
    \item Les indicateurs de dispersion mesurent leur variabilité.
    \item Le graphique doit être choisi selon le type de variable.
    \item Une conclusion statistique doit toujours être interprétée dans son contexte.
\end{itemize}
\end{cadrebleu}

\begin{pagepaysage}

\section{Grand tableau statistique}

\begin{tableausolutionlong}{|p{4cm}|p{8cm}|p{10cm}|}
\hline
\textbf{Indicateur} & \textbf{Description} & \textbf{Interprétation} \\
\hline
Moyenne & Valeur centrale & Tendance générale de la série statistique. \\
\hline
Médiane & Valeur du milieu & Indicateur moins sensible aux valeurs extrêmes. \\
\hline
Variance & Dispersion quadratique & Mesure globale de la variabilité des données. \\
\hline
Écart-type & Racine carrée de la variance & Dispersion moyenne autour de la moyenne. \\
\hline
\end{tableausolutionlong}

\end{pagepaysage}